\begin{abstract}
\noindent
This report describes a university course project: a hands-on study of the Linux
socket API, carried out by building one peer-to-peer chat node in C++23 and giving it
nine interchangeable I/O back ends. Everything above the transport stays fixed: a
shared \code{ChatController}, a single-threaded \code{StoreWorker}, a
\code{SenderWorker} that owns the write policy, and a socket-agnostic
\code{FrameChannel}. The back ends cover blocking thread-per-connection readers,
pooled non-blocking writes, \code{poll(2)} and \code{epoll(7)} reactors over two
sockets per peer (simplex) or one (duplex), reactors with separate read and write
pollers, and one owner thread per peer. A headless test harness and an \code{epoll}
load generator on separate pinned cores drive all nine through fan-in, idle
connections, slow consumers, connection churn, throughput, emulated WAN delay and a
peer cap: 864 runs in total, three repetitions of each point.

The main results:
\begin{itemize}
  \item \textbf{Thread-per-connection collapses at 1,024 peers on two CPUs,} delivering
        10--21\% of the offered load. Below 256 peers it is the cheapest and has the
        lowest tail latency.
  \item \textbf{\code{poll} pays for every idle socket and \code{epoll} does not.}
  \item \textbf{Only a blocking sender lets one slow reader stall every peer.}
  \item \textbf{The second, simplex connection per peer buys no throughput.} It costs
        an extra round trip per handshake, and its blocking dial-back collapses under
        WAN churn.
  \item \textbf{A separate write thread pays off when the single I/O thread is the
        bottleneck.}
  \item \textbf{Under churn, tail latency is set by the shared sender's retry sleep,
        not by the concurrency model.}
\end{itemize}
We close with an untested design hypothesis that combines the measured strengths:
small per-thread peer pools with a balancing policy.
\end{abstract}
